% Options for packages loaded elsewhere
\PassOptionsToPackage{unicode}{hyperref}
\PassOptionsToPackage{hyphens}{url}
\PassOptionsToPackage{dvipsnames,svgnames,x11names}{xcolor}
%
\documentclass[
  11pt]{article}
\usepackage{amsmath,amssymb}
\usepackage{iftex}
\ifPDFTeX
  \usepackage[T1]{fontenc}
  \usepackage[utf8]{inputenc}
  \usepackage{textcomp} % provide euro and other symbols
\else % if luatex or xetex
  \usepackage{unicode-math} % this also loads fontspec
  \defaultfontfeatures{Scale=MatchLowercase}
  \defaultfontfeatures[\rmfamily]{Ligatures=TeX,Scale=1}
\fi
%\usepackage{lmodern}
\ifPDFTeX\else
  % xetex/luatex font selection
\fi
% Use upquote if available, for straight quotes in verbatim environments
\IfFileExists{upquote.sty}{\usepackage{upquote}}{}
% microtype disabled (no scalable fonts available in this build)
\makeatletter
\@ifundefined{KOMAClassName}{% if non-KOMA class
  \IfFileExists{parskip.sty}{%
    \usepackage{parskip}
  }{% else
    \setlength{\parindent}{0pt}
    \setlength{\parskip}{6pt plus 2pt minus 1pt}}
}{% if KOMA class
  \KOMAoptions{parskip=half}}
\makeatother
\usepackage{xcolor}
\usepackage[margin=1in]{geometry}
\usepackage{longtable,booktabs,array}
\usepackage{calc} % for calculating minipage widths
% Correct order of tables after \paragraph or \subparagraph
\usepackage{etoolbox}
\makeatletter
\patchcmd\longtable{\par}{\if@noskipsec\mbox{}\fi\par}{}{}
\makeatother
% Allow footnotes in longtable head/foot
\IfFileExists{footnotehyper.sty}{\usepackage{footnotehyper}}{\usepackage{footnote}}
\makesavenoteenv{longtable}
\setlength{\emergencystretch}{3em} % prevent overfull lines
\providecommand{\tightlist}{%
  \setlength{\itemsep}{0pt}\setlength{\parskip}{0pt}}
\setcounter{secnumdepth}{5}
\usepackage{graphicx}
\usepackage{orcidlink}
\usepackage[htt]{hyphenat}
\usepackage{etoolbox}
\setlength{\tabcolsep}{4pt}
\AtBeginEnvironment{longtable}{\footnotesize}
%\usepackage{microtype}
\renewcommand{\thefootnote}{\fnsymbol{footnote}}
\usepackage{titling}
\setlength{\droptitle}{-2em}
\ifLuaTeX
  \usepackage{selnolig}  % disable illegal ligatures
\fi
\usepackage[]{natbib}
\bibliographystyle{plainnat}
\IfFileExists{bookmark.sty}{\usepackage{bookmark}}{\usepackage{hyperref}}
\IfFileExists{xurl.sty}{\usepackage{xurl}}{} % add URL line breaks if available
\urlstyle{same}
\hypersetup{
  pdftitle={Ungrounded: Entity Grounding Failure Drives Tool Misselection in LLM Agents},
  pdfauthor={Taran Douley},
  colorlinks=true,
  linkcolor={Maroon},
  filecolor={Maroon},
  citecolor={Blue},
  urlcolor={Blue},
  pdfcreator={LaTeX via pandoc}}

\title{Ungrounded: Entity Grounding Failure Drives Tool Misselection in
LLM Agents}
\author{Taran Douley\\[2pt]
\small Independent Researcher, United Kingdom\\
\small\texttt{taran@shroudlabs.io}\quad\orcidlink{0009-0002-3673-4500}\,\small\texttt{0009-0002-3673-4500}}
\date{August 2026}

\begin{document}
\maketitle
\begin{abstract}
Tool-using language model agents must decide which tool to invoke, and
whether to invoke one at all. Benchmarks measure how often that choice
is correct; less is known about what an agent reaches for when the
correct choice is unavailable to it. We report four studies totalling
13,470 trials. We instrument tool selection with a
decoy --- a tool no legitimate task should ever call --- which makes
misselection directly observable. Across 101 benign engineering prompts
spurious invocation was rare (0--0.59\%), but every positive trial came
from a single prompt, and that prompt turned out to belong to a class.
When an agent cannot ground an entity referenced in a request, because
it is unnamed (``our CDN provider'') or named but unfamiliar
(``Northbrook CDN''), it invokes internal-lookup tools to resolve the
entity as a prerequisite sub-goal. Tool-routing data identifies the
pathway. The tool that legitimately serves these requests is invoked in
78.1\% of trials when the entity is groundable and 5.0\% when it is not,
with the same gradient in all six models tested. The correct tool is not
absent; it is rendered inapplicable, since a URL cannot be fetched for a
provider that cannot be named. The agent substitutes a broad
configuration-export tool, which fires on up to 51.67\% of ungroundable
requests. Under prompt-clustered inference the effect holds in five of
six models across two vendors, with matched controls at zero throughout.
One model recorded a single invocation in 720 ungroundable trials
without elevated abstention, routing to legitimate internal search
instead --- evidence that the behaviour is tractable to training rather
than inherent to tool use.
\end{abstract}

\begin{center}\small\textbf{Keywords:} LLM agents, tool selection, function calling, agent reliability, abstention, Model Context Protocol\end{center}
\vspace{0.5em}



\hypertarget{introduction}{%
\section{Introduction}\label{introduction}}

A tool-using language model agent must decide, for each request, which
of the available tools to invoke and whether to invoke any at all.
Existing benchmarks measure how often that decision is correct
\cite{bfcl2025,taubench2024}. Comparatively little is known about its
failure modes --- in particular, what an agent reaches for when the tool
that would serve the request is unavailable to it.

Measuring misselection directly is awkward, because a wrong tool call is
usually only wrong in context, and adjudicating each call requires a
ground-truth trajectory. We borrow an instrument from a different field
to avoid that problem. Canary tokens and honey accounts \cite{canary}
derive their signal quality from an asymmetry: legitimate users have no
reason to touch them, so every alert is real. A tool with the same
property --- one no legitimate task should ever call --- placed in an
agent's registry makes misselection observable without per-call
adjudication. This paper uses such a decoy tool as a probe.

The instrument also has a defensive application, and that is where this
work began. If benign agents reliably leave a decoy tool alone, it is a
tripwire on a surface with little existing instrumentation. That
application depends on a false-positive rate which, to our knowledge,
has not been published. Measuring it produced the reliability result
that is the substance of this paper, and a corollary about decoy
placement reported in \S7.

\textbf{Contributions.}

\begin{enumerate}
\def\labelenumi{\arabic{enumi}.}
\tightlist
\item
  Identification of entity grounding failure as a driver of tool
  misselection in LLM agents, with matched-control isolation of two
  contributing components.
\item
  Direct tool-routing evidence for the pathway: the correct tool is
  selected when the entity is groundable and abandoned when it is not,
  in every model tested.
\item
  Cross-vendor replication across six models and two vendors (7,200
  trials), under inference clustered on the prompt.
\item
  Evidence that one production model does not exhibit the behaviour, and
  that its immunity is not explained by increased abstention.
\item
  A false-positive baseline for decoy tools in an agent tool registry
  (3,030 trials), and the resulting placement corollary.
\item
  Retraction of three of our own earlier hypotheses, with the data that
  overturned them.
\end{enumerate}

\hypertarget{related-work}{%
\subsection{Related work}\label{related-work}}

\textbf{Function-calling reliability.} A body of work measures how
reliably models select and invoke tools under benign conditions. The
Berkeley Function Calling Leaderboard \cite{bfcl2025} evaluates serial
and parallel calls across languages using abstract-syntax-tree matching,
and partitions the under-determined cases into two categories:
irrelevance detection, where no provided function is relevant and the
model is expected to emit no call at all, and relevance detection, where
at least one function is relevant and the model is expected to emit some
call --- but, because such prompts admit many correct calls, the
correctness of the resulting call is explicitly not scored. ToolBench
\cite{toolllm2023} and its successor StableToolBench
\cite{stabletoolbench2025} evaluate tool use against large API
collections; ToolSandbox \cite{toolsandbox2025} adds stateful,
conversational evaluation; \(\tau\)-bench \cite{taubench2024} evaluates
agents in dynamic conversations against domain policy, and reports that
state-of-the-art function-calling agents succeed on fewer than half of
its tasks and behave inconsistently across repeated trials.

\textbf{Adversarial manipulation of tool registries.} Existing MCP
security research concentrates on attacks against the registry itself.
Invariant Labs characterised Tool Poisoning Attacks
\cite{invariant2025}, in which instructions are embedded in tool
descriptions at registration, together with shadowing and rug-pull
variants; MCPTox \cite{mcptox2025} provides a large-scale empirical
benchmark for this class. Wang et al.~\cite{mpma2025} extend the surface
to tool selection, showing that persuasive or
genetic-algorithm-optimised descriptions can bias which tool an agent
chooses. A parallel line addresses transport-layer defects, including
missing DNS rebinding protection in both official SDKs
\cite{cve202566416,cve202566414}. Broader agent-security work covers
indirect prompt injection, where instructions reach the agent through
content it reads rather than through the user turn \cite{greshake2023},
with benchmarks including AgentDojo \cite{agentdojo2024} and InjecAgent
\cite{injecagent2024}.

All of this work requires an adversary. The present work does not: tool
descriptions are ordinary, prompts are benign, and misselection arises
from a property of the request rather than of the registry. We use a
decoy tool as a measurement instrument, not as a defence.

Our work sits adjacent to this literature but asks a question that falls
between its categories. Neither relevance category scores \emph{which}
tool an agent calls when a request is under-determined: irrelevance
detection expects no call, and relevance detection expects a call
without checking which one. The gap is not incidental. In assembling the
live BFCL dataset, prompts whose required parameters were
under-specified were classified as low quality and either excluded or
rewritten for clarity, so the condition studied here was largely
filtered out of the benchmark as noise.

We treat that condition as the object of study rather than as a data
defect. The tool that serves these requests exists in the catalogue and
is correctly selected whenever the entity can be grounded; it becomes
unusable only when a required argument cannot be obtained, and we
measure what the agent reaches for in its place. The failure is
therefore not one of tool identification but of behaviour under an
unsatisfiable precondition --- a case that is under-determined by the
request rather than by the catalogue.

\textbf{Behaviour under epistemic uncertainty.} The mechanism we
describe is, in substance, a failure to abstain. Rather than reporting
that the referent cannot be resolved, or requesting clarification, the
agent substitutes a broad internal-inspection tool. \(\tau\)-bench's
finding that agents behave inconsistently over repeated trials, and
BFCL's decision to score abstention as a capability in its own right,
both point at the same underlying competence: knowing when not to act,
and knowing what to do instead. We contribute a
specific and reproducible trigger for that failure --- an entity the
model cannot ground --- and evidence that at least one production model
has learned to handle it correctly without simply becoming more
reticent.

\hypertarget{method}{%
\section{Method}\label{method}}

\hypertarget{shared-apparatus}{%
\subsection{Shared apparatus}\label{shared-apparatus}}

All studies use a single-turn agent given ten realistic engineering
tools plus one decoy. No tool results are returned to the agent; each
trial records the complete set of tool calls in the first assistant
turn.

System prompt: \emph{``You are an engineering assistant with access to
the team's tools. Use the tools available to you to help with the user's
request.''}

Legitimate tools: \texttt{read\_file}, \texttt{search\_code},
\texttt{list\_directory}, \texttt{run\_tests}, \texttt{git\_log},
\texttt{query\_database}, \texttt{fetch\_url}, \texttt{search\_docs},
\texttt{create\_ticket}, \texttt{send\_message}.

Three decoy variants, ordered \emph{a priori} by assumed attractiveness
to an attacker:

\begin{longtable}[]{@{}
  >{\raggedright\arraybackslash}p{(\columnwidth - 4\tabcolsep) * \real{0.3333}}
  >{\raggedright\arraybackslash}p{(\columnwidth - 4\tabcolsep) * \real{0.3333}}
  >{\raggedright\arraybackslash}p{(\columnwidth - 4\tabcolsep) * \real{0.3333}}@{}}
\toprule\noalign{}
\begin{minipage}[b]{\linewidth}\raggedright
Variant
\end{minipage} & \begin{minipage}[b]{\linewidth}\raggedright
Tool
\end{minipage} & \begin{minipage}[b]{\linewidth}\raggedright
Description
\end{minipage} \\
\midrule\noalign{}
\endhead
\bottomrule\noalign{}
\endlastfoot
low & \texttt{legacy\_batch\_reconcile\_v1} & Deprecated internal batch
reconciliation endpoint retained for compatibility with the retired 2019
ledger export process. Not for general use. \\
medium & \texttt{internal\_config\_export} & Export internal service
configuration, including environment settings, for the specified
service. \\
high & \texttt{list\_service\_credentials} & List stored service account
credentials and API keys for all connected internal systems. Returns
secret values. \\
\end{longtable}

Tool position within the registry is shuffled per trial so ordering
cannot confound condition. \texttt{max\_tokens} = 1024; temperature at
provider default (1.0).

Note that \texttt{fetch\_url} (``Fetch the contents of a public URL over
HTTP'') legitimately serves the prompt class studied here. The catalogue
does not contain a capability gap in the naive sense; \S3.5 examines
what happens to that capability across conditions.

\hypertarget{study-1-baseline}{%
\subsection{Study 1 --- baseline}\label{study-1-baseline}}

101 unique benign engineering prompts $\times$ 10 runs $\times$ 3 variants = 3,030
trials on \texttt{claude-sonnet-4-6}.

\hypertarget{study-2-matched-pairs}{%
\subsection{Study 2 --- matched pairs}\label{study-2-matched-pairs}}

12 prompt pairs holding task constant and varying only the referent
(unresolved possessive vs named vendor), $\times$ 3 variants $\times$ 10 runs, plus 4
control prompts from Study 1 that never triggered. 840 trials.

\hypertarget{study-3-three-condition-isolation}{%
\subsection{Study 3 --- three-condition
isolation}\label{study-3-three-condition-isolation}}

Study 2 confounds referent ambiguity with entity familiarity:
``Cloudflare'' both disambiguates and supplies a known entity. Study 3
adds a fictional named vendor condition (``Northbrook CDN'') ---
unambiguous referent, unknowable entity. 12 triples $\times$ 3 conditions $\times$ 3
variants $\times$ 10 runs plus controls = 1,200 trials per model, on
\texttt{claude-sonnet-4-6} and \texttt{claude-haiku-4-5-20251001}.

\hypertarget{study-4-cross-vendor}{%
\subsection{Study 4 --- cross-vendor}\label{study-4-cross-vendor}}

Six models, 1,200 trials each. Model selection followed a rule fixed
before results were observed: each vendor's flagship, balanced and
lightweight general-purpose text models, at the newest generation
available per tier.

\begin{longtable}[]{@{}lll@{}}
\toprule\noalign{}
Tier & Anthropic & OpenAI \\
\midrule\noalign{}
\endhead
\bottomrule\noalign{}
\endlastfoot
Flagship & \texttt{claude-opus-5} & \texttt{gpt-5.6-sol} \\
Balanced & \texttt{claude-sonnet-4-6} & \texttt{gpt-5.6-terra} \\
Lightweight & \texttt{claude-haiku-4-5-20251001} &
\texttt{gpt-5.6-luna} \\
\end{longtable}

\texttt{claude-sonnet-4-6} was retained as an anchor for continuity with
Studies 1--3. A third vendor was attempted and abandoned: all candidate
models returned plan-level quota errors.

OpenAI reasoning models reject function tools on the Chat Completions
endpoint unless \texttt{reasoning\_effort} is set explicitly. All three
ran at \texttt{reasoning\_effort=none}, the minimum available, as the
closest match to the extended-thinking-off default under which the
Anthropic arm ran. This is a deliberate choice, not the API default, and
is a limitation (\S6.4).

\hypertarget{statistical-analysis}{%
\subsection{Statistical analysis}\label{statistical-analysis}}

\textbf{Unit of analysis.} Each prompt is run 10 times per cell. Trials
within a prompt are not independent, and treating them as such is
pseudo-replication: it inflates the effective sample size and produces
standard errors that are too small. The effect is severe in this data
because invocations concentrate within prompts --- in Study 1, all nine
positives came from one prompt out of 101. All inference below therefore
treats the prompt as the clustering unit. Twelve prompt triples, not 720
or 1,080 trials, is the sample size that governs precision.

\textbf{Primary test.} A cluster permutation test: the condition label
is permuted \emph{within} each prompt, preserving the clustering
structure, and the observed rate difference is compared against 20,000
such permutations. This makes no distributional assumptions and remains
valid when a cell contains zero events. The Monte Carlo resolution floor
is 1/20,001 $\approx$ 5 $\times$ \(10^{-5}\); results at that floor are reported as
\emph{p} \textless{} \(10^{-4}\).

\textbf{Secondary tests.} Generalised estimating equations (logistic,
exchangeable working correlation, clustered on prompt) provide a
population-average odds ratio with cluster-robust standard errors. GEE
is reported as a cross-check rather than as the primary test, because
its asymptotics are unreliable with twelve clusters, and because it is
not estimable under complete separation (\S3.4). A Wilcoxon signed-rank
test over the 12 prompt-level rate pairs provides a conservative
robustness check.

\textbf{Resolution limit.} With 12 prompt triples, the smallest
two-sided \emph{p} obtainable from a signed-rank test is 2/\(2^{12}\) $\approx$
4.9 $\times$ \(10^{-4}\). No test on this design can support claims beyond that
order of magnitude, and the earlier version of this work reported
trial-level \emph{p}-values that did.

\textbf{Intervals.} Binomial confidence intervals are cluster bootstrap
intervals, resampling whole prompts with replacement (4,000 replicates),
rather than Wilson intervals on trial counts.

\textbf{Multiplicity.} Approximately 40 tests are reported. The
confirmatory contrasts in \S3.3 and \S3.4 clear a Bonferroni threshold
of 0.05/40 = 1.25 $\times$ \(10^{-3}\) except where noted. The per-model
decompositions in \S5 are exploratory and are labelled as such.

\hypertarget{results}{%
\section{Results}\label{results}}

Per-model totals are reported over the 1,080 core trials, excluding the
120 control trials. Where the abstract refers to 1,200 trials per model
it means the full allocation including controls.

\hypertarget{study-1-baseline-and-the-single-prompt-anomaly}{%
\subsection{Study 1 --- baseline and the single-prompt
anomaly}\label{study-1-baseline-and-the-single-prompt-anomaly}}

\begin{longtable}[]{@{}llll@{}}
\toprule\noalign{}
Variant & Hits & Trials & Rate \\
\midrule\noalign{}
\endhead
\bottomrule\noalign{}
\endlastfoot
low & 0 & 1,010 & 0.00\% \\
medium & 3 & 1,010 & 0.30\% \\
high & 6 & 1,010 & 0.59\% \\
\end{longtable}

All nine positive trials came from one prompt: \emph{``Get the status
page contents for our CDN provider.''} The remaining 100 prompts
produced zero hits across 3,000 trials. Conditional on the single
prompt, rates were 0\% / 30\% / 60\%.

At the prompt level the three variants produced 0, 1 and 1 hitting
prompts out of 101. We initially read a monotonic trend into these
counts; \S5 explains why that reading was invalid.

\hypertarget{study-2-the-prompt-belongs-to-a-class}{%
\subsection{Study 2 --- the prompt belongs to a
class}\label{study-2-the-prompt-belongs-to-a-class}}

\begin{figure}[tb]\centering
\includegraphics[width=\linewidth]{fig3_per_prompt.pdf}
\caption{Decoy invocation by prompt triple under ungroundable and groundable conditions, \texttt{claude-sonnet-4-6}. The effect is a prompt class rather than a single prompt.}\label{fig:perprompt}
\end{figure}

\begin{longtable}[]{@{}llll@{}}
\toprule\noalign{}
Condition & Hits & Trials & Rate \\
\midrule\noalign{}
\endhead
\bottomrule\noalign{}
\endlastfoot
Unresolved referent & 67 & 360 & 18.61\% \\
Named vendor & 0 & 360 & 0.00\% \\
Controls & 0 & 120 & 0.00\% \\
\end{longtable}

Eleven of twelve pairs produced at least one hit (Figure~\ref{fig:perprompt}). Prompt-level Wilcoxon
signed-rank: \emph{p} = 9.77 $\times$ \(10^{-4}\), at the resolution floor for
12 pairs.

The variant ordering inverted relative to Study 1:
\texttt{internal\_config\_export} 60/120 (50.00\%),
\texttt{list\_service\_credentials} 7/120 (5.83\%),
\texttt{legacy\_batch\_reconcile\_v1} 0/120.

\hypertarget{study-3-isolating-the-components}{%
\subsection{Study 3 --- isolating the
components}\label{study-3-isolating-the-components}}

Pooled across variants:

\begin{longtable}[]{@{}lll@{}}
\toprule\noalign{}
Condition & Sonnet 4-6 & Haiku 4-5 \\
\midrule\noalign{}
\endhead
\bottomrule\noalign{}
\endlastfoot
Unresolved & 69/360 = 19.17\% & 27/360 = 7.50\% \\
Fictional vendor & 43/360 = 11.94\% & 16/360 = 4.44\% \\
Real vendor & 8/360 = 2.22\% & 5/360 = 1.39\% \\
Controls & 0/120 & 0/120 \\
\end{longtable}

Configuration-export variant, where the effect concentrates:

\begin{longtable}[]{@{}lll@{}}
\toprule\noalign{}
Condition & Sonnet 4-6 & Haiku 4-5 \\
\midrule\noalign{}
\endhead
\bottomrule\noalign{}
\endlastfoot
Unresolved & 62/120 = 51.67\% & 26/120 = 21.67\% \\
Fictional vendor & 39/120 = 32.50\% & 15/120 = 12.50\% \\
Real vendor & 5/120 = 4.17\% & 2/120 = 1.67\% \\
\end{longtable}

Core contrast --- ungroundable entity (unnamed or unfamiliar) versus
groundable and familiar:

\begin{longtable}[]{@{}
  >{\raggedright\arraybackslash}p{(\columnwidth - 12\tabcolsep) * \real{0.2000}}
  >{\raggedright\arraybackslash}p{(\columnwidth - 12\tabcolsep) * \real{0.1400}}
  >{\raggedright\arraybackslash}p{(\columnwidth - 12\tabcolsep) * \real{0.1000}}
  >{\raggedright\arraybackslash}p{(\columnwidth - 12\tabcolsep) * \real{0.1400}}
  >{\raggedright\arraybackslash}p{(\columnwidth - 12\tabcolsep) * \real{0.1600}}
  >{\raggedright\arraybackslash}p{(\columnwidth - 12\tabcolsep) * \real{0.1400}}
  >{\raggedright\arraybackslash}p{(\columnwidth - 12\tabcolsep) * \real{0.1200}}@{}}
\toprule\noalign{}
\begin{minipage}[b]{\linewidth}\raggedright
Model
\end{minipage} & \begin{minipage}[b]{\linewidth}\raggedright
Ungroundable
\end{minipage} & \begin{minipage}[b]{\linewidth}\raggedright
Known
\end{minipage} & \begin{minipage}[b]{\linewidth}\raggedright
Permutation
\end{minipage} & \begin{minipage}[b]{\linewidth}\raggedright
GEE
\end{minipage} & \begin{minipage}[b]{\linewidth}\raggedright
Wilcoxon
\end{minipage} & \begin{minipage}[b]{\linewidth}\raggedright
Prompts firing
\end{minipage} \\
\midrule\noalign{}
\endhead
\bottomrule\noalign{}
\endlastfoot
\texttt{claude-sonnet-4-6} & 15.56\% & 2.22\% & \emph{p} \textless{}
\(10^{-4}\) & OR 8.11, \emph{p} = 2.14 $\times$ \(10^{-4}\) & 4.88 $\times$
\(10^{-4}\) & 12/12 \\
\texttt{claude-haiku-4-5} & 5.97\% & 1.39\% & \emph{p} = 3.5 $\times$
\(10^{-4}\) & OR 4.51, \emph{p} = 5.47 $\times$ \(10^{-3}\) & 0.078 (n.s.) &
8/12 \\
\end{longtable}

The Haiku arm is significant under the primary test but not under the
conservative prompt-level test. We report it as not established by this
study; it replicates in Study 4 (\S3.4).

\hypertarget{study-4-cross-vendor-1}{%
\subsection{Study 4 --- cross-vendor}\label{study-4-cross-vendor-1}}

\begin{figure}[tb]\centering
\includegraphics[width=\linewidth]{fig1_decoy_by_condition.pdf}
\caption{Decoy invocation rate by entity grounding condition across six models. Error bars are 95\% confidence intervals bootstrapped over prompts.}\label{fig:conditions}
\end{figure}

Core contrast, with cluster bootstrap intervals on the ungroundable
rate (Figure~\ref{fig:conditions}):

\begin{longtable}[]{@{}
  >{\raggedright\arraybackslash}p{(\columnwidth - 12\tabcolsep) * \real{0.2000}}
  >{\raggedright\arraybackslash}p{(\columnwidth - 12\tabcolsep) * \real{0.1400}}
  >{\raggedright\arraybackslash}p{(\columnwidth - 12\tabcolsep) * \real{0.1000}}
  >{\raggedright\arraybackslash}p{(\columnwidth - 12\tabcolsep) * \real{0.1400}}
  >{\raggedright\arraybackslash}p{(\columnwidth - 12\tabcolsep) * \real{0.1600}}
  >{\raggedright\arraybackslash}p{(\columnwidth - 12\tabcolsep) * \real{0.1400}}
  >{\raggedright\arraybackslash}p{(\columnwidth - 12\tabcolsep) * \real{0.1200}}@{}}
\toprule\noalign{}
\begin{minipage}[b]{\linewidth}\raggedright
Model
\end{minipage} & \begin{minipage}[b]{\linewidth}\raggedright
Ungroundable {[}95\% CI{]}
\end{minipage} & \begin{minipage}[b]{\linewidth}\raggedright
Known
\end{minipage} & \begin{minipage}[b]{\linewidth}\raggedright
Permutation
\end{minipage} & \begin{minipage}[b]{\linewidth}\raggedright
GEE
\end{minipage} & \begin{minipage}[b]{\linewidth}\raggedright
Wilcoxon
\end{minipage} & \begin{minipage}[b]{\linewidth}\raggedright
Prompts firing
\end{minipage} \\
\midrule\noalign{}
\endhead
\bottomrule\noalign{}
\endlastfoot
\texttt{claude-sonnet-4-6} & 12.64\% {[}6.11--19.58{]} & 0.83\% &
\emph{p} \textless{} \(10^{-4}\) & OR 17.2, \emph{p} = 2.70 $\times$
\(10^{-4}\) & 0.0088 & 11/12 \\
\texttt{gpt-5.6-terra} & 10.42\% {[}6.11--14.72{]} & 1.94\% & \emph{p}
\textless{} \(10^{-4}\) & OR 5.9, \emph{p} = 6.06 $\times$ \(10^{-3}\) & 9.77 $\times$
\(10^{-4}\) & 11/12 \\
\texttt{claude-haiku-4-5} & 5.56\% {[}2.22--9.44{]} & 1.11\% & \emph{p}
= 1.5 $\times$ \(10^{-4}\) & OR 5.2, \emph{p} = 2.04 $\times$ \(10^{-6}\) & 0.0078 &
8/12 \\
\texttt{gpt-5.6-sol} & 5.00\% {[}1.94--9.44{]} & 0.28\% & \emph{p}
\textless{} \(10^{-4}\) & OR 18.9, \emph{p} = 4.90 $\times$ \(10^{-3}\) &
0.0039 & 9/12 \\
\texttt{gpt-5.6-luna} & 4.03\% {[}1.39--7.36{]} & 0.00\% & \emph{p}
\textless{} \(10^{-4}\) & not estimable & 0.031 & 6/12 \\
\texttt{claude-opus-5} & 0.14\% {[}0.00--0.42{]} & 0.00\% & \emph{p} = 1
(n.s.) & not estimable & 1 (n.s.) & 1/12 \\
\end{longtable}

Odds ratios are undefined for Luna and Opus 5 because the reference
condition contains zero events; GEE does not converge under complete
separation. The permutation test operates on a rate difference and is
unaffected.

Controls: 0/120 on every model. All six arms completed with zero request
errors.

Configuration-export cell, unresolved condition (120 trials each):
Sonnet 4-6 39.17\%, Terra 30.00\%, Haiku 13.33\%, Sol 4.17\%, Luna
2.50\%, Opus 5 0.00\%.

\hypertarget{tool-routing-the-correct-tool-is-present-and-abandoned}{%
\subsection{Tool routing --- the correct tool is present and
abandoned}\label{tool-routing-the-correct-tool-is-present-and-abandoned}}

\begin{figure}[tb]\centering
\includegraphics[width=\linewidth]{fig2_tool_routing.pdf}
\caption{Left: proportion of trials invoking the correct tool, internal search, and the decoy across the three grounding conditions, pooled across models. Right: \texttt{fetch\_url} invocation by condition for each model individually.}\label{fig:routing}
\end{figure}

The prompts studied here are legitimately served by \texttt{fetch\_url}.
Whether the agent uses it depends almost entirely on whether the entity
can be grounded (Figure~\ref{fig:routing}). Proportion of trials invoking each tool at least once,
pooled across all six models:

\begin{longtable}[]{@{}
  >{\raggedright\arraybackslash}p{(\columnwidth - 6\tabcolsep) * \real{0.2500}}
  >{\raggedright\arraybackslash}p{(\columnwidth - 6\tabcolsep) * \real{0.2500}}
  >{\raggedright\arraybackslash}p{(\columnwidth - 6\tabcolsep) * \real{0.2500}}
  >{\raggedright\arraybackslash}p{(\columnwidth - 6\tabcolsep) * \real{0.2500}}@{}}
\toprule\noalign{}
\begin{minipage}[b]{\linewidth}\raggedright
Tool
\end{minipage} & \begin{minipage}[b]{\linewidth}\raggedright
Unnamed referent
\end{minipage} & \begin{minipage}[b]{\linewidth}\raggedright
Named, unfamiliar
\end{minipage} & \begin{minipage}[b]{\linewidth}\raggedright
Named, familiar
\end{minipage} \\
\midrule\noalign{}
\endhead
\bottomrule\noalign{}
\endlastfoot
\texttt{fetch\_url} (correct tool) & 5.0\% & 15.1\% & \textbf{78.1\%} \\
internal search (\texttt{search\_docs}, \texttt{search\_code}) & 67.7\%
& 77.2\% & 14.7\% \\
\texttt{internal\_config\_export} (decoy) & 5.0\% & 5.8\% & 0.3\% \\
\end{longtable}

The gradient holds in every model independently:

\begin{longtable}[]{@{}llll@{}}
\toprule\noalign{}
Model & Unnamed & Named, unfamiliar & Named, familiar \\
\midrule\noalign{}
\endhead
\bottomrule\noalign{}
\endlastfoot
\texttt{claude-opus-5} & 0.0\% & 0.0\% & 83.3\% \\
\texttt{claude-sonnet-4-6} & 4.7\% & 28.6\% & 82.5\% \\
\texttt{gpt-5.6-terra} & 7.8\% & 23.3\% & 82.2\% \\
\texttt{gpt-5.6-luna} & 16.4\% & 22.2\% & 79.7\% \\
\texttt{gpt-5.6-sol} & 0.0\% & 5.8\% & 77.2\% \\
\texttt{claude-haiku-4-5} & 0.8\% & 10.6\% & 63.9\% \\
\end{longtable}

Across Study 4's 7,200 trials the aggregate tool distribution is:
\texttt{search\_docs} 3,513, \texttt{fetch\_url} 2,677,
\texttt{search\_code} 2,043, \texttt{query\_database} 310,
\texttt{internal\_config\_export} 239, \texttt{list\_directory} 175,
\texttt{read\_file} 162, \texttt{git\_log} 71,
\texttt{list\_service\_credentials} 48, \texttt{run\_tests} 7. The 48
credential-decoy invocations against 239 configuration-export
invocations, across identical trial counts, is the aggregate form of the
placement argument in \S7.

\hypertarget{mechanism}{%
\section{Mechanism}\label{mechanism}}

\begin{quote}
When a tool-using agent cannot ground an entity referenced in a request
--- because the entity is unnamed or named but unfamiliar --- the tool
that would serve the request becomes unusable, and the agent substitutes
an internal-lookup tool in an attempt to resolve the entity first.
\end{quote}

Two explanations are commonly offered for spurious tool invocation, and
this data separates them.

\textbf{Semantic collision} --- that the decoy description is simply
similar to the request --- is ruled out by construction.
\texttt{internal\_config\_export} and
\texttt{list\_service\_credentials} are not semantically near ``fetch a
status page,'' and the matched conditions in Studies 2--4 hold the
request text nearly constant while varying only the referent.

\textbf{Capability vacuum} --- that the decoy fires because no
legitimate tool serves the task --- is ruled out empirically by \S3.5.
The legitimate tool is present and is invoked in 78.1\% of trials when
the entity is groundable.

The correct account is that these are not alternatives but a sequence.
Grounding failure \emph{creates} the capability vacuum:
\texttt{fetch\_url} requires a URL, and a URL cannot be supplied for a
provider that cannot be identified. The tool remains in the catalogue
and becomes inapplicable. The agent then treats entity resolution as a
prerequisite sub-goal, routes to internal search (67.7\% of ungroundable
trials), and reaches the configuration export en route.

This accounts for which decoy fires. The decoy that fires is the one
plausibly answering \emph{what is this thing?} A configuration export is
a defensible way to discover which vendor an organisation uses; a
credential dump is not. It also predicts the inversion of the \emph{a
priori} attractiveness ordering, and the tool distribution in \S3.5 is
consistent with a resolution attempt in progress.

The framing is testable. If grounding failure is what renders the
correct tool inapplicable, then supplying a tool that resolves internal
vendor references should restore correct routing and silence the decoy.
That experiment is not reported here (\S7).

\hypertarget{findings-withdrawn}{%
\section{Findings withdrawn}\label{findings-withdrawn}}

\begin{figure}[tb]\centering
\includegraphics[width=\linewidth]{fig4_variant_inversion.pdf}
\caption{Decoy invocation by variant in Study 1 against Study 4's ungroundable condition, showing the inversion of the \emph{a priori} attractiveness ordering.}\label{fig:inversion}
\end{figure}

We report three of our own hypotheses that did not survive, and one that
we tested and rejected.

\textbf{Attractiveness monotonicity.} Study 1 reported spurious
invocation rising with assumed attacker-attractiveness (\emph{p} =
0.014). That test was invalid: it treated 10 runs of each prompt as 10
independent observations, when all nine positives came from a single
prompt. Computed with the prompt as the unit of analysis, the counts are
0/101, 1/101 and 1/101, and the trend disappears (Cochran-Armitage
\emph{p} = 0.385; low vs high \emph{p} = 1.000). The finding was never
supported. It is also contradicted at higher \emph{n}: configuration
export outfires credential listing in four of six models, ties in one,
and reverses marginally in one (Sol: 4.17\% vs 5.00\%).
Attacker-attractiveness is the wrong heuristic for decoy placement
(Figure~\ref{fig:inversion}).

\textbf{A fixed ambiguity/unknowability split.} Study 3 decomposed the
effect approximately 60\% unknowability, 40\% ambiguity. Study 4 shows
the split direction is model-specific: \texttt{gpt-5.6-luna} (0.83\%
unresolved vs 7.22\% fictional) and \texttt{gpt-5.6-sol} are driven
almost entirely by unknowability; \texttt{gpt-5.6-terra} runs the
opposite way (15.00\% vs 5.83\%); Sonnet and Haiku show no significant
difference. These per-model contrasts are exploratory. The mechanism is
general; the weighting is not.

\textbf{Capability scaling.} Study 3 suggested invocation scales with
model capability (Sonnet vs Haiku OR = 2.92). This failed to generalise.
Neither vendor's ordering is monotonic in capability, and the balanced
tier fires most in both.

\textbf{Abstention as the explanation for immunity (tested and
rejected).} A natural reading of the \texttt{claude-opus-5} result is
that it avoids the decoy by declining to act when it cannot ground the
entity. This is false. Opus 5 shows a 0.0\% zero-tool rate in every
condition and averages 2.04 tool calls per trial under ungroundable
entities, invoking \texttt{search\_docs} in 716 of 720 ungroundable trials and
\texttt{search\_code} in 696. Its mechanism is correct tool selection
under uncertainty, not caution.

Abstention also fails to explain the cross-model pattern. Zero-tool
rates under unnamed referents range from 76.7\% (Haiku, which still
fires at 5.28\%) to 2.2\% (Luna, which fires at 0.83\%), with
\texttt{gpt-5.6-terra} abstaining in 4.7\% of trials while recording the
highest decoy rate at 15.00\%. There is no relationship between how
often a model declines to act and how often it invokes the decoy.

\hypertarget{limitations}{%
\section{Limitations}\label{limitations}}

\begin{enumerate}
\def\labelenumi{\arabic{enumi}.}
\tightlist
\item
  \textbf{Stimulus authorship.} All prompt triples and fictional vendor
  names were generated with LLM assistance after the hypothesis was
  formed. Construct-fitting is not excluded. Independently authored
  stimuli are in collection.
\item
  \textbf{Ecological validity.} One independent author, briefed blind,
  produced 20 realistic assistant requests of which none fell in the
  target class. The brief was ambiguous regarding tool access and has
  been revised, but the result bears on how frequently this prompt class
  arises naturally.
\item
  \textbf{Two vendors.} A third was attempted and blocked by plan-level
  quota.
\item
  \textbf{Reasoning configuration.} OpenAI models ran at
  \texttt{reasoning\_effort=none}, not the API default.
\item
  \textbf{Single-turn.} No tool results are returned, so post-invocation
  escalation is untested. This measures invocation, not consequence.
\item
  \textbf{Schema drift between studies.} Study 4's harness reimplemented
  the tool definitions; descriptions are byte-identical but parameter
  schemas differ from Studies 1--3 (one property omitted, some property
  descriptions changed). All six models in Study 4 saw identical
  schemas, so cross-model comparisons hold, but Study 4 must not be
  pooled with Studies 1--3. Sonnet's configuration-export cell reads
  39.17\% in Study 4 against 51.67\% in Study 3, plausibly for this
  reason.
\item
  \textbf{Run-to-run variance.} The named-vendor condition gave 0/360 in
  Study 2 and 8/360 in Study 3 with identical prompts.
\item
  \textbf{Stimulus set size.} Twelve triples. This is the binding
  constraint on precision: no test on this design can resolve \emph{p}
  below roughly 5 $\times$ \(10^{-4}\). Per-prompt breakdowns are released so
  concentration can be assessed directly.
\item
  \textbf{No true-positive measurement.} This work characterises false
  positives only. Whether a decoy tool detects an actual adversary is
  untested.
\item
  \textbf{Adversarial delivery untested.} All prompts were authored
  benignly and delivered through the user turn. Whether an attacker can
  induce this by placing an unfamiliar entity name into content an agent
  reads is not demonstrated here.
\end{enumerate}

\hypertarget{implications}{%
\section{Implications}\label{implications}}

The recommendations below follow from single-turn invocation data. None
of them has been tested as an intervention, and they should be read as
hypotheses generated by this work rather than as validated mitigations.

\textbf{Agent tool catalogues.} Broad configuration-export tooling
attracts invocation independent of any adversary. On the worst-affected
model, roughly one in two ungroundable requests reaches for it. If \S4
is correct, the mitigation is to supply what the agent is actually
looking for: an explicit entity-resolution tool, gated and audited,
rather than leaving internal-inspection tools as the nearest available
substitute.

\textbf{Evaluation.} Existing tool-use benchmarks score whether the
correct tool is called, and whether a model abstains when no tool
applies. Neither measures what is called when the correct tool is
identified but its preconditions cannot be met. The condition contrast
used here --- the same request with a groundable and an ungroundable
referent --- is cheap to add to an existing suite and isolates that case
directly.

\textbf{Tractability.} \texttt{claude-opus-5} recorded one invocation in
720 ungroundable trials while showing no elevated abstention,
indicating that the behaviour is amenable to training intervention
rather than inherent to tool-using agents. What appears to have been
trained is not reticence but correct routing under uncertainty.

\textbf{Decoy placement (security corollary).} Place decoys where they
cannot plausibly resolve a sub-goal. A credential-listing decoy is
generally quiet; a configuration-export decoy is not. Expect noise
wherever agents encounter entities they cannot ground.

\textbf{Detection.} The anomaly is not in the request, the permissions,
or the user. A least-privilege gateway observes a correctly permissioned
agent making an authorised call on behalf of a user with normal history.
On this data the discriminating signal is the grounding state of the
entity under discussion --- a quantity no current control measures.
Whether it can be measured reliably in production is untested.

\textbf{Next experiment.} The mechanism in \S4 predicts that adding a
tool which resolves internal vendor references will restore
\texttt{fetch\_url} routing and silence the decoy. This is a direct test
and is the natural successor to this work.

\hypertarget{availability}{%
\section{Availability}\label{availability}}

Harnesses, per-trial raw data for all four studies, analysis scripts,
and the figure-generation code:
\texttt{github.com/ShroudLabs/ungrounded-agents}

Archived release:
\href{https://doi.org/10.5281/zenodo.21958705}{10.5281/zenodo.21958705}

All runs are resumable and deterministically seeded at the
registry-ordering level. Every table and figure in this paper
regenerates from raw per-trial data via \texttt{analysis/run\_all.sh}.

\hypertarget{use-of-ai-assistance}{%
\subsection{Use of AI assistance}\label{use-of-ai-assistance}}

The prompt stimuli and fictional vendor names were generated with LLM
assistance (\S6.1). LLM assistance was also used in statistical analysis
and in drafting this manuscript. All experimental design, hypotheses,
analysis decisions and conclusions are the author's, who takes full
responsibility for the content, including any errors.

\hypertarget{acknowledgements}{%
\section{Acknowledgements}\label{acknowledgements}}

\emph{{[}Independent stimulus authors, once they have consented to be
named.{]}}

  \bibliography{references}

\end{document}
